\section{Discussion}

The main results suggest that TAMR-DTI is most useful when the prediction task
benefits from pair-specific representation learning. Rather than treating the
drug and protein as two fixed encodings that are fused only at the end, the
model lets the protein context influence conformer aggregation and lets the
drug context condition protein encoding before final interaction prediction.
This design matches the structure of DTI prediction: the same ligand can expose
different useful geometric views for different targets, and the same protein
sequence can carry different interaction evidence depending on the ligand. The
benchmark results therefore support the value of target-aware modulation, while
the threshold-dependent metrics remind us that ranking quality and calibrated
decision behavior are not the same problem.

The conformer-count analysis should be read in this light. It does not imply
that adding more conformers will always improve performance. Conformers are
computationally generated geometric hypotheses, and later conformer slots may
include less relevant or noisier geometries. The useful point is narrower:
binding prediction can benefit from exposing multiple plausible molecular
geometries, provided that the model has a target-conditioned mechanism for
weighting them. This is why the conformer branch is better viewed as a
protein-guided geometry aggregation module than as a simple multi-conformer
augmentation trick.

The protein-side refinement modules play a complementary role. FiLM modulation
injects ligand information into protein feature extraction, while the Mamba
block refines the ordered protein-token stream before cross-modal matching.
TAMR-DTI deliberately keeps attention-based interaction modeling in the final
fusion stage, because DTI prediction still requires explicit drug-protein
alignment. Mamba is therefore used as a controlled residual refinement on the
protein side, not as a wholesale replacement for cross-modal attention. The
ablation results are the main evidence for this design choice. The
population-level interpretability statistics on all \InterpNumSamples{} BioSNAP
test pairs and the DHFR case study in Section~\ref{fig:case_study_dhfr} provide
two complementary checks on this picture: at the population level the gated
refinement is small and stable across prediction groups (mean
$R_t=\InterpMambaRefineMean$, learned gate \InterpMambaGate), and on at least
one well-characterized target the residue-level refinement aligns with
UniProt-annotated binding residues (DHFR, fold \CaseStudyDHFRFold,
$p_{\mathrm{Stouffer}}=\CaseStudyDHFRPComb$, \CaseStudyDHFRNSig{} of
\CaseStudyDHFRNRows{} TP rows individually significant). The remaining
curated targets in our case study do not pass the same test, so the
visualizations should still be read as model-internal evidence and as a
positive example for DHFR, rather than as a general residue-level
binding-site predictor.

Several limitations remain. The reported baselines come from the literature and
do not provide a fully matched paired-testing protocol. The ablation and
conformer-count studies are single-seed diagnostics on BioSNAP, so they support
mechanistic interpretation but should not be treated as formal significance
tests. The generated conformers are not experimentally resolved binding poses,
and the current random-split benchmarks do not fully test scaffold-level or
target-family generalization. The case study is also limited: only
\CaseStudyNumSigTargets{} of the \CaseStudyNumTargets{} curated UniProt
targets passes the Stouffer-combined Fisher's-exact test, and the protein-side
Mamba branch is trained on interaction labels alone without residue-level
supervision. Future work should therefore use matched multi-seed ablations,
harder out-of-distribution splits, calibration-aware threshold selection, and
stronger structural supervision when reliable binding pose or binding-site
information is available.
